\section{\textbf{Referências}}
%\vspace*{-24 pt}
\begin{SingleSpace}
	\noindent BARAFFE, I. et al. New evolutionary models for pre-main sequence and main sequence low-mass stars down to the hydrogen-burning limit.\textbf{ Astronomy \& Astrophysics}, v. 577, p. A42, April 2015
	\vspace*{12 pt}
	
	\noindent DAHM, S. E. Reexamining the lithium depletion boundary in the pleiades and the inferred
	age of the cluster. \textbf{The Astrophysical Journal}, v. 813, n. 2, p. 108, 2015.	\vspace*{12 pt}
	
	
	\noindent DELFOSSE, X. et al. Accurate masses of very low mass stars. iv. improved mass- luminosity relations.\textbf{ Astronomy and Astrophysics}, v. 364, p. 217–224, December 2000.
		\vspace*{12 pt}
	
	
	\noindent Gaia Collaboration et al. Gaia data release 3. summary of the content and survey
	properties. \textbf{Astronomy and Astrophysics}, v. 674, p. A1, 2023.	\vspace*{12 pt}
	
	
	\noindent	 GOSSAGE, S. et al. Age determinations of the hyades, praesepe, and pleiades via mesa models with rotation. \textbf{The Astrophysical Journal}, v. 863, n. 1, p. 67, 2018.	\vspace*{12 pt}
	
	
	\noindent  LODIEU, N. et al. A 5d view of the alpha per, pleiades, and praesepe clusters. \textbf{Astronomy \& Astrophysics}, v. 628, p. A66, 2019.
	\vspace*{12 pt}
	
	\noindent SIESS, L.; DUFOUR, E.; FORESTINI, M. An internet server for update pre-main sequence tracks of low- and intermediate-mass stars. \textbf{Astronomy \& Astrophysics}, 2000.	\vspace*{1 pt}


	\noindent STASSUN, K. G. et al. The revised tess input catalog and candidate target list. \textbf{The Astronomical Journal}, v. 158, p. 138, October 2019.	\vspace*{12 pt}
	
	
\end{SingleSpace}

\section{\textbf{Publicações}}
\noindent XVII Semana da Física da Uesc em 20 de Outubro de 2022;\\
\noindent XVIII Semana da Física da Uesc em 6 de Outubro de 2023.
